% =====================================================================
% A2 — Section 2: Background.  DRAFT 2, 3 Oct 2026.
% Restructured after the A1 feedback: purpose before machinery, concept
% before maths, one job per paragraph, Fig. 1 carries the argument.
% Self-contained (A2 must be readable without A1); notation follows A1.
% Draft 2 cuts the displayed effective loss: only H is needed downstream.
% =====================================================================

\section{Background}
\label{sec:background}


Liu et al.~\cite{liu2022grokking} study grokking on a task whose correct
solution is known in advance, so that one can check directly whether a
network has found it. Their argument has two stages summarised in
Fig.~\ref{fig:pipeline}. First, the training data must contain enough information
to determine the solution, and second, the training must then actually find it.

\begin{figure}[t]
  \centering
  % =====================================================================
% A2 — Figure 1: the two-stage argument ("data determine" -> "training
% finds"). Original methodological figure; drawn in TikZ so it stays vector
% and uses the document font. Fits one column (column width 8.6 cm):
% stage labels in a left column (x < -1.6), flow in the middle and right.
% =====================================================================
\begin{tikzpicture}[
    font=\footnotesize,
    box/.style={draw, rounded corners=2pt, align=center, inner sep=2.5pt,
                minimum height=0.7cm, fill=white},
    q/.style={draw, diamond, aspect=2.2, align=center, inner sep=0.5pt,
              fill=white},
    arr/.style={-{Stealth[length=4pt]}, thick},
    stage/.style={font=\scriptsize, text=black!70, align=left,
                  text width=1.95cm, anchor=west},
  ]
  % --- flow -------------------------------------------------------------
  \node[box, text width=3.5cm] (D) at (0.5, 0)
       {Training set $D$, fraction $r$};
  \node[box, text width=3.7cm] (A) at (0.5, -1.05)
       {Equal-sum pairs $\to$ equations\\ $E_i + E_j = E_m + E_n$};
  \node[q] (Q) at (0.5, -2.25) {$\lambda_3 > 0$\,?};
  \node[box, text width=1.55cm, fill=black!6] (N) at (3.25, -2.25)
       {Not\\ determined};

  \node[box, text width=3.9cm] (T) at (0.5, -3.75)
       {Representation learning\\ \emph{vs.}\ decoder overfitting};
  \node[box, text width=2.15cm] (S) at (-0.75, -5.0)
       {Structure first:\\ RQI $\to 1$};
  \node[box, text width=2.15cm, fill=black!6] (M) at (1.85, -5.0)
       {Decoder first:\\ memorisation};

  \draw[arr] (D) -- (A);
  \draw[arr] (A) -- (Q);
  \draw[arr] (Q) -- node[above, font=\scriptsize] {no} (N);
  \draw[arr] (Q) -- node[pos=0.22, right, font=\scriptsize] {yes} (T);
  \draw[arr] (T.south -| S.north) -- (S.north);
  \draw[arr] (T.south -| M.north) -- (M.north);

  % --- stage divider and labels ------------------------------------------
  \draw[densely dashed, black!40] (-3.55, -2.95) -- (4.15, -2.95);
  \node[stage] at (-3.55, -1.6)
       {\textbf{Do the data}\\ \textbf{determine} the\\ representation?\\ (C1)};
  \node[stage] at (-3.55, -4.2)
       {\textbf{Does training}\\ \textbf{find} it?\\ (C2, C3)};
\end{tikzpicture}

  \caption{The argument tested in this study. Upper part, whether the
  training set determines the linear representation depends only on the
  equal-sum pairs it contains, summarised by $\lambda_3 > 0$ (claim C1).
  Lower part, whether training then finds that representation depends on
  the competition between representation learning and decoder overfitting
  (claims C2 and C3).}
  \label{fig:pipeline}
\end{figure}

\subsection{Task and model}
\label{sec:task}

The task needs a known correct solution, so Liu et al.\ use non-modular
addition on the symbols $\{0, \dots, p-1\}$, here with $p = 10$. Because addition is a commutative operation, $(a,b)$ and $(b,a)$ are treated as
one sample, so the full dataset $D_0$ contains $|D_0| = p(p+1)/2 = 55$ pairs.
A training set $D \subset D_0$ is drawn at random, the remaining pairs are
held out, and $r = |D|/|D_0|$ is the \emph{training data fraction}.

The model gives each symbol a trainable embedding
$\mathbf{E}_a \in \mathbb{R}^{d_{\mathrm{in}}}$, adds the two embeddings of a
pair, and passes the sum to a decoder network:
%
\begin{equation}
  (a,b) \;\longmapsto\; \mathrm{Dec}\!\left(\mathbf{E}_a + \mathbf{E}_b\right).
  \label{eq:toymodel}
\end{equation}
%
Each possible sum $c$ has a fixed random target vector $\mathbf{Y}_c$, and the
model predicts the sum of those target is closest to its output. Since the
addition of embeddings is built into the model, the only way it can answer an
unseen pair is through the arrangement of its embeddings. As in the paper's Fig.~4, I use one-dimensional embeddings ($d_{\mathrm{in}} = 1$) and write
$\mathbf{R} = [\mathbf{E}_0, \dots, \mathbf{E}_{p-1}]$ for the representation.

\subsection{Representation quality}
\label{sec:rqi}

To measure whether the embeddings have the right arrangement, Liu et al.\ use
pairs with equal sums. Say $i + j = m + n$ and the model maps both pairs to the
same point, $\mathbf{E}_i + \mathbf{E}_j = \mathbf{E}_m + \mathbf{E}_n$. Then
the decoder cannot tell the two pairs apart, so a correct answer for one
transfers to the other. Such a quadruple is called a \emph{parallelogram}.

For a model with zero training loss and an injective decoder, every two
training pairs with the same sum must form a parallelogram
(\cite{liu2022grokking}, Prop.~2). These are
%
\begin{equation}
  P_0(D) = \left\{ (i,j,m,n) \,\middle|\,
    \begin{array}{l}
      (i,j) \in D,\ (m,n) \in D, \\[-1pt]
      i+j = m+n
    \end{array} \right\}.
  \label{eq:permissible}
\end{equation}
%
The \emph{representation quality index} is the fraction of all parallelograms
of the full dataset, $P_0 \equiv P_0(D_0)$, that the representation realises
within a tolerance $\delta$. Writing $P(\mathbf{R}, \delta)$ for the
parallelograms in $P_0$ with
$|\mathbf{E}_i + \mathbf{E}_j - \mathbf{E}_m - \mathbf{E}_n| \le \delta$,
%
\begin{equation}
  \mathrm{RQI}(\mathbf{R}) = \frac{|P(\mathbf{R}, \delta)|}{|P_0|} \in [0, 1].
  \label{eq:rqi}
\end{equation}
%
Evenly spaced embeddings, $\mathbf{E}_k = \mathbf{a} + k\mathbf{b}$, realise
every parallelogram and give $\mathrm{RQI} = 1$. This is the \emph{linear representation}, where random embeddings give $\mathrm{RQI} \approx 0$.

\subsection{Effective theory and the grokking rate}
\label{sec:efftheory}

The effective theory is about whether a training set contains enough
parallelograms to force the linear representation and, it is, then how quickly
training gets there. To make this tractable, Liu et al.\ leave out the decoder, so
the normalised embeddings descend a loss on violated parallelograms.

In one dimension each parallelogram in $P_0(D)$ is a linear equation
$E_i + E_j = E_m + E_n$, and these equations form the set $A(P_0(D))$. The effective loss $\ell_{\mathrm{eff}} = \ell_0/Z_0$ is the mean
squared violation $\ell_0$ of the equations, normalised by
$Z_0 = \sum_k |\tilde{\mathbf{E}}_k|^2$. And, because $\ell_0$ is quadratic,
gradient descent on it is linear, $d\mathbf{R}/dt = -\mathbf{H}\mathbf{R}$,
with the Hessian
%
\begin{equation}
  \mathbf{H}_{ij} = \frac{1}{Z_0}
    \frac{\partial^2 \ell_0}{\partial \mathbf{E}_i \partial \mathbf{E}_j}.
  \label{eq:hessian}
\end{equation}
%
Each eigenvector of $\mathbf{H}$ is a direction in which the embeddings can
move, and its eigenvalue, sorted as $\lambda_1 \le \dots \le \lambda_p$, is
the rate at which deviations in that direction die out. A zero eigenvalue means
that the training equations do not constrain that direction.

Translating or scaling all embeddings would not break a parallelogram, so
$\lambda_1 = \lambda_2 = 0$ for every training set. If
$\lambda_3 = 0$, a non-linear representation fits the training equations equally
well. If $\lambda_3 > 0$, every other direction is constrained, the nullity
$n_0$ of $A(P_0(D))$ is two, and the linear representation is unique up to
these two changes.

In that case, $\lambda_3$ is the slowest rate that matters, so it sets the timescale for
approaching the linear representation. Liu et al.\ call it the
\emph{grokking rate} and give the time towards the linear structure as
%
\begin{equation}
  t_h = \frac{1}{\lambda_3}, \qquad n_h = \frac{1}{\lambda_3\, \eta},
  \label{eq:grokrate}
\end{equation}
%
where $\eta$ is the embedding learning rate and $n_h$ the number of steps.

\subsection{Claims under test}
\label{sec:claims}

I test three claims of~\cite{liu2022grokking} for $p = 10$:
%
\begin{itemize}
  \item[\textbf{C1}] \emph{(\cite{liu2022grokking}, Fig.~4a.)} The probability that a random training
    set gives $\lambda_3 > 0$ rises sharply at a critical data fraction
    $r_c \approx 0.4$.
  \item[\textbf{C2}] \emph{(\cite{liu2022grokking}, Fig.~4b.)} ``The number of steps to reach
    $\mathrm{RQI} > 0.95$ is seen to have a phase transition at $r_c = 0.4$,
    agreeing with the proposed effective theory.''
  \item[\textbf{C3}] \emph{(\cite{liu2022grokking}, Fig.~5b.)} Among trained networks, the steps to
    $\mathrm{RQI} > 0.95$ decrease with $\lambda_3$, as predicted by
    $t \sim 1/\lambda_3$.
\end{itemize}
%
Claim 1 (C1) depends only on the training set, the upper part of
Fig.~\ref{fig:pipeline}. Claim 2 (C2) and Claim 3 (C3) depend on what training does with it, the
lower part.
